\documentclass[aspectratio=169,10pt,xcolor={dvipsnames,table}]{beamer}

% Tema Moderno Metropolis
\usetheme{metropolis}

% Pacotes Fundamentais
\usepackage[utf8]{inputenc}
\usepackage[T1]{fontenc}
\usepackage[brazilian]{babel}
\usepackage{graphicx}
\usepackage{booktabs}
\usepackage{amsmath}
\usepackage{amssymb}
\usepackage{fontawesome5}
\usepackage{tcolorbox}
\tcbuselibrary{skins,raster}
\usepackage{bookmark}

% Cores Institucionais e Tecnologicas (IFCE + Tech Modern)
\definecolor{IFCEGreen}{RGB}{47, 158, 65}       % Verde IFCE oficial
\definecolor{IFCEDarkGreen}{RGB}{24, 107, 43}    % Verde escuro institucional
\definecolor{IFCERed}{RGB}{200, 16, 46}          % Vermelho IFCE oficial
\definecolor{TechDark}{RGB}{30, 41, 59}          % Azul ardosia escuro para texto
\definecolor{TechLight}{RGB}{248, 250, 252}      % Cinza suave para fundos
\definecolor{TechBorder}{RGB}{226, 232, 240}     % Borda sutil
\definecolor{AccentBlue}{RGB}{2, 132, 199}       % Azul vibrante (web/nuvem)
\definecolor{AccentPurple}{RGB}{124, 58, 237}    % Roxo IA/Inovacao
\definecolor{AccentOrange}{RGB}{234, 88, 12}     % Laranja Empreendedorismo

% Configuracao de Cores no Metropolis Beamer
\setbeamercolor{palette primary}{fg=white, bg=IFCEDarkGreen}
\setbeamercolor{palette secondary}{fg=white, bg=IFCEGreen}
\setbeamercolor{palette tertiary}{fg=white, bg=TechDark}
\setbeamercolor{structure}{fg=IFCEDarkGreen}
\setbeamercolor{titlelike}{fg=IFCEDarkGreen}
\setbeamercolor{background canvas}{bg=white}
\setbeamercolor{normal text}{fg=TechDark}
\setbeamercolor{progress bar}{fg=IFCERed, bg=IFCEGreen!30}
\setbeamercolor{alerted text}{fg=IFCERed}

% Configuracao do cabecalho e rodape do Metropolis
\metroset{
  progressbar=frametitle,
  block=fill,
  numbering=fraction
}

% Boxes e Cards Estilizados com TColorBox
\newtcolorbox{moderncard}[2][]{
  enhanced,
  colback=white,
  colframe=TechDark!75!black,
  colbacktitle=TechDark!80!black,
  coltitle=white,
  fonttitle=\bfseries\footnotesize,
  arc=2mm,
  boxrule=0.8pt,
  left=2mm, right=2mm, top=1.5mm, bottom=1.5mm,
  title={#2},
  #1
}

\newtcolorbox{greencard}[2][]{
  enhanced,
  colback=white,
  colframe=IFCEDarkGreen,
  colbacktitle=IFCEDarkGreen,
  coltitle=white,
  fonttitle=\bfseries\footnotesize,
  arc=2mm,
  boxrule=0.8pt,
  left=2mm, right=2mm, top=1.5mm, bottom=1.5mm,
  title={#2},
  #1
}

\newtcolorbox{bluecard}[2][]{
  enhanced,
  colback=white,
  colframe=AccentBlue!90!black,
  colbacktitle=AccentBlue!90!black,
  coltitle=white,
  fonttitle=\bfseries\footnotesize,
  arc=2mm,
  boxrule=0.8pt,
  left=2mm, right=2mm, top=1.5mm, bottom=1.5mm,
  title={#2},
  #1
}

\newtcolorbox{purplecard}[2][]{
  enhanced,
  colback=white,
  colframe=AccentPurple!85!black,
  colbacktitle=AccentPurple!85!black,
  coltitle=white,
  fonttitle=\bfseries\footnotesize,
  arc=2mm,
  boxrule=0.8pt,
  left=2mm, right=2mm, top=1.5mm, bottom=1.5mm,
  title={#2},
  #1
}

\newtcolorbox{redcard}[2][]{
  enhanced,
  colback=white,
  colframe=IFCERed,
  colbacktitle=IFCERed,
  coltitle=white,
  fonttitle=\bfseries\footnotesize,
  arc=2mm,
  boxrule=0.8pt,
  left=2mm, right=2mm, top=1.5mm, bottom=1.5mm,
  title={#2},
  #1
}

\newtcolorbox{statcard}[2][]{
  enhanced,
  colback=TechLight,
  colframe=TechBorder,
  arc=2mm,
  boxrule=0.6pt,
  left=2mm, right=2mm, top=1.5mm, bottom=1.5mm,
  halign=center,
  #1
}

% Metadados da Apresentacao
\title[ADS - IFCE Campus Tauá]{\textbf{Tecnologia em Análise e Desenvolvimento de Sistemas}}
\subtitle{Construa sua Carreira Tech Global a Partir do Sertão dos Inhamuns}
\author[IFCE Tauá]{Instituto Federal de Educação, Ciência e Tecnologia do Ceará\texorpdfstring{\\ \textbf{Campus Tauá}}{ - Campus Tauá}}
\institute[IFCE]{Curso Superior de Graduação Tecnológica}
\date{\today}

\begin{document}

% Secao 1: Capa e Visao Geral do Curso
% Secao 1: Capa e Visao Geral do Curso

% Slide 1: Capa Customizada de Alto Impacto Visual
\begin{frame}[plain]
  \begin{columns}[c]
    \begin{column}{0.38\textwidth}
      \centering
      \includegraphics[width=0.78\textwidth]{logo_ifce_taua.png}
      \vspace{0.3cm}
      
      {\scriptsize \textbf{\textcolor{IFCEDarkGreen}{MINISTÉRIO DA EDUCAÇÃO}}} \\
      \vspace{0.05cm}
      {\tiny \textbf{Secretaria de Educação Profissional e Tecnológica}}
    \end{column}
    
    \begin{column}{0.60\textwidth}
      \vspace{-0.1cm}
      {\scriptsize \textcolor{IFCERed}{\textbf{\faGraduationCap\ CURSO SUPERIOR DE GRADUAÇÃO TECNOLÓGICA}}}
      \vspace{0.12cm}
      
      {\large \textbf{\textcolor{IFCEDarkGreen}{Análise e Desenvolvimento\\de Sistemas (ADS)}}}
      
      \vspace{0.15cm}
      {\footnotesize \textbf{Construa sua Carreira Tech Global a Partir dos Inhamuns}}
      
      \vspace{0.25cm}
      \begin{tcolorbox}[
        enhanced,
        colback=IFCEGreen!8,
        colframe=IFCEGreen!60,
        arc=2mm,
        boxrule=0.6pt,
        left=2.5mm, right=2.5mm, top=1.8mm, bottom=1.8mm
      ]
        \scriptsize
        \faRocket\ \textbf{Duração:} Apenas 2,5 anos (5 semestres letivos) \\
        \faLaptopCode\ \textbf{Formação:} 2.100h alinhadas ao mercado moderno \\
        \faAward\ \textbf{Diploma Federal:} 100\% Gratuito e Reconhecido pelo MEC
      \end{tcolorbox}
      
      \vspace{0.15cm}
      {\tiny \textcolor{TechDark!70}{Tauá -- CE \quad | \quad Projeto Pedagógico do Curso (PPC)}}
    \end{column}
  \end{columns}
\end{frame}

% Slide 2: Roteiro da Apresentação
\begin{frame}{Roteiro da Apresentação}{\faListUl\ O que você vai descobrir sobre o curso de ADS}
  \begin{columns}[t]
    \begin{column}{0.48\textwidth}
      \begin{moderncard}{\faCompass\ Fundamentos \& Carreira}
        \begin{itemize}\setlength{\itemsep}{0.2em}\footnotesize
          \item \textbf{Visão Geral:} Duração, turnos e estrutura ágil.
          \item \textbf{Mercado \& Trabalho Remoto:} Oportunidades globais sem sair dos Inhamuns.
          \item \textbf{Grade Curricular da SBC:} Os 4 eixos formativos essenciais.
          \item \textbf{Carreira Dev, IA \& LLMs:} Como a matriz prepara você para a fronteira tech.
        \end{itemize}
      \end{moderncard}
    \end{column}
    
    \begin{column}{0.48\textwidth}
      \begin{moderncard}{\faGlobe\ Vivência \& Permanência}
        \begin{itemize}\setlength{\itemsep}{0.2em}\footnotesize
          \item \textbf{Extensão \& Prática Real:} PIM I e II e impacto na sociedade local.
          \item \textbf{Além da Sala de Aula:} Intercâmbio Internacional, Pesquisa e Monitoria.
          \item \textbf{Infraestrutura de Ponta:} Laboratórios climatizados e biblioteca digital.
          \item \textbf{Ingresso \& Assistência:} 35 vagas via SISU e suporte financeiro ao estudante.
        \end{itemize}
      \end{moderncard}
    \end{column}
  \end{columns}
\end{frame}

% Slide 3: Visão Geral do Curso
\begin{frame}{Visão Geral do Curso}{\faClock\ Formação Superior Ágil, Estratégica e Gratuita}
  \begin{columns}[t]
    \begin{column}{0.5\textwidth}
      \begin{greencard}{\faCheckCircle\ Estrutura Acadêmica}
        \begin{itemize}\setlength{\itemsep}{0.25em}\footnotesize
          \item \textbf{Duração:} 2,5 anos (5 semestres letivos).
          \item \textbf{Turnos:} Presencial nos turnos Matutino ou Vespertino.
          \item \textbf{Titulação:} Tecnólogo em Análise e Desenvolvimento de Sistemas.
          \item \textbf{Carga Horária Total:} \textbf{2.100 horas} intensivas:
            \begin{itemize}\tiny
              \item 1.840h de disciplinas obrigatórias
              \item 160h de disciplinas optativas diversificadas
              \item 100h de atividades complementares
              \item 210h (10\%) para Curricularização da Extensão
            \end{itemize}
        \end{itemize}
      \end{greencard}
    \end{column}
    
    \begin{column}{0.47\textwidth}
      \begin{bluecard}{\faGraduationCap\ O Diferencial do Tecnólogo}
        \begin{itemize}\setlength{\itemsep}{0.25em}\footnotesize
          \item \textbf{Diploma de Nível Superior:} Mesma validade legal que bacharelados para concursos e pós-graduação.
          \item \textbf{Foco na Empregabilidade:} Matriz alinhada com as demandas reais da indústria de software.
          \item \textbf{Velocidade de Retorno:} Em 2 anos e meio você já é um graduado com portfólio prático.
          \item \textbf{Continuidade Acadêmica:} Acesso direto a Especializações, Mestrados e Doutorados.
        \end{itemize}
      \end{bluecard}
    \end{column}
  \end{columns}
\end{frame}


% Secao 2: O Profissional de ADS e as Fronteiras do Mercado
% Secao 2: O Profissional de ADS e as Fronteiras do Mercado

\begin{frame}{O Profissional de ADS}{\faUserTie\ O que faz um Tecnólogo em Análise e Desenvolvimento de Sistemas?}
  \begin{columns}[t]
    \begin{column}{0.48\textwidth}
      \begin{moderncard}{\faLaptopCode\ Do Conceito à Produção}
        \scriptsize
        O tecnólogo domina todo o ciclo de vida do software:
        \begin{itemize}\setlength{\itemsep}{0.2em}
          \item \textbf{Elicitação e Análise:} Entende a dor do cliente e traduz necessidades de negócio em requisitos técnicos.
          \item \textbf{Arquitetura e Projeto:} Modela bancos de dados, diagramas e arquiteturas escaláveis.
          \item \textbf{Codificação Full-Stack:} Constrói soluções para Web, dispositivos móveis e nuvem.
          \item \textbf{Qualidade e Testes:} Aplica testes automatizados, auditoria e boas práticas de código limpo.
        \end{itemize}
      \end{moderncard}
    \end{column}
    
    \begin{column}{0.48\textwidth}
      \begin{greencard}{\faBriefcase\ Carreiras e Funções em Alta}
        \scriptsize
        Amplas oportunidades no ecossistema tech:
        \begin{itemize}\setlength{\itemsep}{0.2em}
          \item \textbf{Desenvolvedor Full-Stack / Mobile}
          \item \textbf{Engenheiro de Software / Arquiteto}
          \item \textbf{Engenheiro de Dados \& Analista de BI}
          \item \textbf{Especialista em DevOps \& Cloud}
          \item \textbf{Engenheiro de Inteligência Artificial / Prompt Engineer}
          \item \textbf{Gerente de Projetos de TI / Scrum Master}
          \item \textbf{Consultor e Fundador de Startups}
        \end{itemize}
      \end{greencard}
    \end{column}
  \end{columns}
\end{frame}

\begin{frame}{Mercado \& Região dos Inhamuns}{\faGlobeAmericas\ Trabalho Remoto Global e Desenvolvimento Local}
  \begin{columns}[t]
    \begin{column}{0.48\textwidth}
      \begin{bluecard}{\faHome\ A Revolução do Trabalho Remoto}
        \begin{itemize}\setlength{\itemsep}{0.25em}\footnotesize
          \item \textbf{Salários Globais em Tauá:} Trabalhe para capitais brasileiras ou empresas internacionais recebendo em reais, dólar ou euro sem sair de casa.
          \item \textbf{Fim do Êxodo Forçado:} Construa uma carreira de prestígio internacional perto da sua família e da sua região.
          \item \textbf{Alto Poder de Compra:} Renda de tecnologia com o custo de vida atrativo do interior cearense.
        \end{itemize}
      \end{bluecard}
    \end{column}
    
    \begin{column}{0.48\textwidth}
      \begin{greencard}{\faSeedling\ Impacto no Sertão dos Inhamuns}
        \begin{itemize}\setlength{\itemsep}{0.25em}\footnotesize
          \item \textbf{Transformação Digital Regional:} Demanda em alta por softwares no comércio, saúde, agronegócio e gestão pública local.
          \item \textbf{Ecossistema Empreendedor:} Capacidade para fundar sua própria software house ou startup regional.
          \item \textbf{Fixação de Talentos:} Retenção de mentes brilhantes para gerar riqueza e inovação em Tauá.
        \end{itemize}
      \end{greencard}
    \end{column}
  \end{columns}
\end{frame}


% Secao 3: A Grade Curricular e os Eixos da SBC
% Secao 3: A Grade Curricular e os Eixos da SBC

\begin{frame}{Estrutura Curricular Conforme a SBC}{\faCubes\ Formação de Excelência Alinhada às Diretrizes Nacionais da SBC e ENADE}
  \begin{columns}[t]
    \begin{column}{0.48\textwidth}
      \begin{bluecard}{\faProjectDiagram\ 1. Análise e Projeto de Software (360h)}
        \tiny
        \textbf{Competência:} Modelar e conceber sistemas para problemas reais.
        \begin{itemize}\setlength{\itemsep}{0.05em}
          \item Engenharia de Software \& Gestão de Projetos Ágeis
          \item Análise e Projeto de Sistemas (UML / Arquitetura)
          \item Interação Humano-Computador (Design UI/UX)
          \item Padrões de Projeto de Software (Design Patterns GoF)
        \end{itemize}
      \end{bluecard}
      
      \vspace{0.1cm}
      
      \begin{greencard}{\faServer\ 3. Infraestrutura para Software (400h)}
        \tiny
        \textbf{Competência:} Ambientes tecnológicos, redes e computação em nuvem.
        \begin{itemize}\setlength{\itemsep}{0.05em}
          \item Redes de Computadores \& Segurança da Informação
          \item Sistemas Operacionais \& Administração Linux
          \item Desenvolvimento e Operações (\textbf{DevOps \& CI/CD})
        \end{itemize}
      \end{greencard}
    \end{column}
    
    \begin{column}{0.48\textwidth}
      \begin{purplecard}{\faCode\ 2. Implementação de Software (1.360h)}
        \tiny
        \textbf{Competência:} Programação avançada, bancos de dados e IA.
        \begin{itemize}\setlength{\itemsep}{0.05em}
          \item Lógica, POO e Estruturas de Dados Avançadas
          \item Bancos de Dados Relacionais (SQL) e Não-Relacionais (\textbf{NoSQL})
          \item Full-Stack Web (Web I e II) \& Apps Mobile nativos/híbridos
          \item \textbf{Inteligência Computacional \& Ciência de Dados}
          \item Testes Automatizados e Qualidade de Software
        \end{itemize}
      \end{purplecard}
      
      \vspace{0.1cm}
      
      \begin{moderncard}{\faUserFriends\ 4. Empreendedorismo \& Pessoal (480h)}
        \tiny
        \textbf{Competência:} Liderança, visão de negócio, comunicação e ética.
        \begin{itemize}\setlength{\itemsep}{0.05em}
          \item Empreendedorismo, Inovação e Modelos de Startups
          \item Inglês Técnico, Comunicação e Expressão
          \item Ética, Responsabilidade Socioambiental e Projeto Social
        \end{itemize}
      \end{moderncard}
    \end{column}
  \end{columns}
\end{frame}

\begin{frame}{Evolução Semestral da Formação}{\faStream\ Um Percurso Passo a Passo do Básico ao Nível Avançado}
  \begin{columns}[t]
    \begin{column}{0.48\textwidth}
      \begin{greencard}{\faSeedling\ Ciclo de Fundamentação (1º ao 3º Sem.)}
        \scriptsize
        \textbf{Construindo as Bases do Desenvolvedor:}
        \begin{itemize}\setlength{\itemsep}{0.2em}
          \item \textbf{1º Semestre (400h):} Introdução à Programação, Banco de Dados, Tecnologias Web, Pensamento Computacional, Fundamentos de Matemática e Inglês Técnico.
          \item \textbf{2º Semestre (400h):} POO, Estrutura de Dados, Engenharia de Software, Redes de Computadores e Empreendedorismo.
          \item \textbf{3º Semestre (400h):} Programação Web I, Dispositivos Móveis, Bancos NoSQL, Análise e Projeto de Sistemas, Sistemas Operacionais e IHC.
        \end{itemize}
      \end{greencard}
    \end{column}
    
    \begin{column}{0.48\textwidth}
      \begin{bluecard}{\faRocket\ Ciclo Avançado e Produção (4º e 5º Sem.)}
        \scriptsize
        \textbf{Inteligência, Nuvem e Prática Real:}
        \begin{itemize}\setlength{\itemsep}{0.2em}
          \item \textbf{4º Semestre (400h):} \textbf{Inteligência Computacional}, Programação Web II, Padrões de Projeto, Segurança da Informação e \textbf{PIM I} (Projeto Integrador).
          \item \textbf{5º Semestre (400h):} \textbf{DevOps}, Testes e Qualidade de Software, \textbf{PIM II}, Projeto Social e Optativas Avançadas.
          \item \textbf{Optativas de Destaque:} Ciência de Dados, Jogos Digitais, Tópicos em IA e Métodos Ágeis.
        \end{itemize}
      \end{bluecard}
    \end{column}
  \end{columns}
  
  \vspace{0.15cm}
  \begin{center}
    \scriptsize \textcolor{IFCEDarkGreen}{\faCheckCircle\ \textbf{Resultado:} Uma progressão pedagógica pensada para você construir portfólio desde o primeiro ano.}
  \end{center}
\end{frame}


% Secao 4: Impulso na Carreira Dev, Inteligência Artificial (LLMs) e Pesquisa
% Secao 4: Impulso na Carreira Dev, Inteligencia Artificial (LLMs) e Pesquisa

\begin{frame}{Alavancando a Carreira de Desenvolvedor}{\faRocket\ Da Lógica Básica à Engenharia de Software de Alto Desempenho}
  \begin{columns}[t]
    \begin{column}{0.32\textwidth}
      \begin{bluecard}{\faCode\ Full-Stack \& Mobile}
        \tiny
        \textbf{O Coração do Mercado}
        \begin{itemize}\setlength{\itemsep}{0.15em}
          \item \textbf{Web I e II:} APIs RESTful, front-ends reativos e back-ends escaláveis.
          \item \textbf{Mobile:} Apps para smartphones Android e ecossistemas móveis.
          \item \textbf{SQL + NoSQL:} Persistência relacional e bancos de alta escala.
        \end{itemize}
      \end{bluecard}
    \end{column}
    
    \begin{column}{0.32\textwidth}
      \begin{purplecard}{\faDraftingCompass\ Arquitetura \& Testes}
        \tiny
        \textbf{O Diferencial Sênior}
        \begin{itemize}\setlength{\itemsep}{0.15em}
          \item \textbf{Padrões de Projeto:} Código limpo e desacoplado com padrões GoF.
          \item \textbf{Testes e Qualidade:} Automação de testes unitários, integração e TDD.
          \item \textbf{Eng. de Software:} Requisitos e arquitetura sustentável.
        \end{itemize}
      \end{purplecard}
    \end{column}
    
    \begin{column}{0.32\textwidth}
      \begin{greencard}{\faCloud\ DevOps \& Nuvem}
        \tiny
        \textbf{Entrega Contínua}
        \begin{itemize}\setlength{\itemsep}{0.15em}
          \item \textbf{DevOps:} Pipelines de integração e entrega contínua (CI/CD).
          \item \textbf{Docker \& Nuvem:} Implantação e orquestração ágil de contêineres.
          \item \textbf{Segurança:} Práticas de DevSecOps e proteção de dados.
        \end{itemize}
      \end{greencard}
    \end{column}
  \end{columns}
  
  \vspace{0.15cm}
  \begin{statcard}
    \scriptsize \textcolor{TechDark}{\faAward\ \textbf{Vantagem Competitiva:} Formação completa que habilita o egresso a disputar vagas júnior/pleno em qualquer empresa do Brasil ou do exterior.}
  \end{statcard}
\end{frame}

\begin{frame}{A Nova Fronteira: IA, LLMs e Pesquisa Científica}{\faBrain\ Aliando Modelos de Linguagem e Inovação Acadêmica}
  \begin{columns}[t]
    \begin{column}{0.48\textwidth}
      \begin{purplecard}{\faRobot\ Inteligência Artificial \& LLMs}
        \scriptsize
        Na disciplina de \textbf{Inteligência Computacional (80h)} e optativas de IA e Dados:
        \begin{itemize}\setlength{\itemsep}{0.15em}
          \item \textbf{Fundamentos Sólidos:} Machine learning, redes neurais e representação de conhecimento vetorial.
          \item \textbf{Copilotos \& Agentes:} Engenharia de prompt e integração de APIs de LLMs para automação inteligente.
          \item \textbf{Arquiteturas RAG:} Criação de sistemas de busca semântica unindo bancos de dados e LLMs.
          \item \textbf{Auditoria Crítica:} Combate a alucinações com rigor de engenharia de software e testes.
        \end{itemize}
      \end{purplecard}
    \end{column}
    
    \begin{column}{0.48\textwidth}
      \begin{greencard}{\faMicroscope\ Crescimento na Pesquisa Acadêmica}
        \scriptsize
        Portas abertas para Mestrado e Inovação:
        \begin{itemize}\setlength{\itemsep}{0.15em}
          \item \textbf{Bolsas de Pesquisa (PIBIC / PIBITI):} Projetos remunerados orientados por mestres e doutores.
          \item \textbf{IA para Demandas Locais:} Aplicações no agronegócio, recursos hídricos, saúde e educação no semiárido.
          \item \textbf{Publicações Científicas:} Apresentação de artigos em congressos da SBC (ERBASE, CBIE, IHC).
          \item \textbf{Ponte para Pós-Graduação:} Vantagem curricular decisiva para ingresso direto em Mestrados e Doutorados.
        \end{itemize}
      \end{greencard}
    \end{column}
  \end{columns}
\end{frame}


% Secao 5: Diferenciais Práticos e Curricularização da Extensão
% Secao 5: Diferenciais Praticos e Curricularizacao da Extensao

\begin{frame}{Diferenciais Práticos: PIM I \& II e PPS}{\faTools\ Da Teoria à Prática: Desenvolvendo Softwares Reais na Graduação}
  \begin{columns}[t]
    \begin{column}{0.48\textwidth}
      \begin{bluecard}{\faLayerGroup\ Projetos Integradores (PIM I e II)}
        \begin{itemize}\setlength{\itemsep}{0.35em}\footnotesize
          \item \textbf{Integração Multidisciplinar:} Componentes obrigatórios no 4º e 5º semestres com 160h combinadas.
          \item \textbf{Criação de Produto Completo:} As equipes concebem, projetam, codificam, testam e implantam um software funcional.
          \item \textbf{Portfólio para Recrutadores:} O estudante se forma com sistemas reais publicados no GitHub e em produção para apresentar em entrevistas de emprego.
        \end{itemize}
      \end{bluecard}
    \end{column}
    
    \begin{column}{0.48\textwidth}
      \begin{greencard}{\faUserCheck\ Prática Profissional Supervisionada (PPS)}
        \begin{itemize}\setlength{\itemsep}{0.35em}\footnotesize
          \item \textbf{40 Horas de Vivência de Mercado:} Distribuídas diretamente dentro de PIM I e PIM II (20h em cada).
          \item \textbf{Ambiente Simulado de Software House:} Metodologias ágeis, reuniões diárias (dailies), revisões de sprint e entrega de valor.
          \item \textbf{Supervisão Docente Direta:} Feedback contínuo de professores mestres e doutores com experiência na indústria.
        \end{itemize}
      \end{greencard}
    \end{column}
  \end{columns}
\end{frame}

\begin{frame}{Curricularização da Extensão: Impacto Social}{\faHandsHelping\ 10\% do Curso Dedicado à Transformação da Comunidade Local}
  \begin{columns}[t]
    \begin{column}{0.48\textwidth}
      \begin{greencard}{\faHandHoldingHeart\ O que é a Extensão Curricularizada?}
        \begin{itemize}\setlength{\itemsep}{0.35em}\footnotesize
          \item \textbf{210 Horas Obrigatórias (10\% da CH Total):} Integradas progressivamente do 2º ao 5º semestre.
          \item \textbf{Diálogo com a Sociedade:} Ações gratuitas e diretas que conectam a sala de aula às dores reais da população.
          \item \textbf{Disciplinas Envolvidas:} Empreendedorismo, Ética, Dispositivos Móveis, PIM I e II, e Projeto Social.
        \end{itemize}
      \end{greencard}
    \end{column}
    
    \begin{column}{0.48\textwidth}
      \begin{bluecard}{\faCity\ Impacto nos Sertões dos Inhamuns}
        \begin{itemize}\setlength{\itemsep}{0.35em}\footnotesize
          \item \textbf{Soluções para o Comércio \& Produtores:} Criação de sistemas de gestão, apps e canais digitais para arranjos produtivos locais.
          \item \textbf{Inclusão Digital \& Cidadania:} Oficinas gratuitas de capacitação digital para jovens, escolas públicas e agricultores.
          \item \textbf{Desenvolvimento de Soft Skills:} O discente aprimora empatia, oratória, escuta ativa, liderança e compromisso ético-social.
        \end{itemize}
      \end{bluecard}
    \end{column}
  \end{columns}
\end{frame}


% Secao 6: Oportunidades Além da Sala de Aula (Intercâmbio, Bolsas e Monitoria)
% Secao 6: Oportunidades Alem da Sala de Aula

\begin{frame}{Intercâmbio Internacional IFCE}{\faPlaneDeparture\ Estude um Semestre no Exterior com Bolsa Integral!}
  \begin{columns}[t]
    \begin{column}{0.52\textwidth}
      \begin{bluecard}{\faGlobe\ O Mundo ao Alcance do Estudante de Tauá}
        \scriptsize
        \begin{itemize}\setlength{\itemsep}{0.2em}
          \item \textbf{Tradição desde 2016:} Anualmente, o IFCE Campus Tauá seleciona discentes para estudar 1 semestre em universidades no exterior.
          \item \textbf{Custeio Completo:} Bolsa mensal cobrindo passagens aéreas, alojamento, alimentação e seguro-saúde internacional.
          \item \textbf{Auxílio Pré-Embarque:} Recurso institucional para emissão de passaporte e taxas de visto consular.
          \item \textbf{Aproveitamento Curricular:} As disciplinas cursadas fora são integralizadas no seu curso no IFCE.
        \end{itemize}
      \end{bluecard}
    \end{column}
    
    \begin{column}{0.44\textwidth}
      \begin{statcard}
        \scriptsize \textcolor{IFCEDarkGreen}{\textbf{\faAward\ Requisitos Acessíveis}}
        \vspace{0.1cm}
        \begin{itemize}\setlength{\itemsep}{0.15em}\tiny
          \item Estar entre o 2º e o 4º semestre do curso.
          \item Coeficiente acadêmico (IRA) $\ge$ 7,0.
          \item Proficiência no idioma do país de destino.
          \item Histórico acadêmico e carta de motivação.
        \end{itemize}
      \end{statcard}
      \vspace{0.1cm}
      \begin{tcolorbox}[enhanced, colback=IFCERed!5, colframe=IFCERed, arc=2mm, boxrule=0.6pt, left=2mm, right=2mm, top=1.5mm, bottom=1.5mm]
        \tiny \faPassport\ \textbf{Experiência Global:} Um diferencial de peso no currículo para competir em posições tech internacionais.
      \end{tcolorbox}
    \end{column}
  \end{columns}
\end{frame}

\begin{frame}{Bolsas, Iniciação Científica e Monitoria}{\faMedal\ Apoio Financeiro e Aceleração do Conhecimento Acadêmico}
  \begin{columns}[t]
    \begin{column}{0.48\textwidth}
      \begin{greencard}{\faFlask\ Iniciação Científica \& Tecnológica (PIBIC/PIBITI)}
        \scriptsize
        \begin{itemize}\setlength{\itemsep}{0.25em}
          \item \textbf{Bolsas Mensais de Pesquisa:} Financiadas pelo IFCE, CNPq e FUNCAP.
          \item \textbf{Mentoria Exclusiva:} Trabalho direto lado a lado com professores doutores e mestres.
          \item \textbf{Inovação e Propriedade:} Registro de softwares, patentes e artigos científicos.
        \end{itemize}
      \end{greencard}
    \end{column}
    
    \begin{column}{0.48\textwidth}
      \begin{purplecard}{\faChalkboardTeacher\ Programas de Monitoria Acadêmica}
        \scriptsize
        \begin{itemize}\setlength{\itemsep}{0.25em}
          \item \textbf{Monitorias Remuneradas e Voluntárias:} Apoio aos colegas em Programação, Dados e Redes.
          \item \textbf{Fixação Profunda:} Ensinar consolida o domínio técnico e oratória profissional.
          \item \textbf{Pontuação Curricular:} Horas validadas e vantagem em seleções de mestrado e empregos.
        \end{itemize}
      \end{purplecard}
    \end{column}
  \end{columns}
\end{frame}


% Secao 7: Infraestrutura de Ponta (Laboratórios e Biblioteca)
% Secao 7: Infraestrutura de Ponta

\begin{frame}{Infraestrutura de Ponta: Laboratórios}{\faMicrochip\ Ambientes Climatizados e Equipamentos Dedicados à Prática Intensiva}
  \begin{columns}[t]
    \begin{column}{0.48\textwidth}
      \begin{bluecard}{\faDesktop\ Laboratórios de Informática 1 e 2}
        \scriptsize
        \begin{itemize}\setlength{\itemsep}{0.22em}
          \item \textbf{Mais de 60 Computadores:} Máquinas dedicadas exclusivamente às aulas práticas e ao desenvolvimento de projetos dos discentes.
          \item \textbf{Conforto Térmico:} Ambientes totalmente climatizados para estudo e prática contínua.
          \item \textbf{Acessibilidade Completa:} Postos de trabalho adaptados para Pessoas com Deficiência (PcD).
          \item \textbf{Ambientes Prontos:} Linux, ferramentas de programação, IDEs, Docker e compiladores configurados.
        \end{itemize}
      \end{bluecard}
    \end{column}
    
    \begin{column}{0.48\textwidth}
      \begin{greencard}{\faNetworkWired\ Laboratório de Redes \& Telecom}
        \scriptsize
        \begin{itemize}\setlength{\itemsep}{0.18em}
          \item \textbf{Racks Industriais:} Switches gerenciáveis de alta performance e Patch Panels Cat 5E.
          \item \textbf{Equipamentos de Ponta:} Roteadores Gigabit Dual-Band AC1750 e centrais Gateway IP com ramais VoIP.
          \item \textbf{Prática de Redes na Realidade:} Montagem de cabos, roteamento avançado, firewalls e servidores.
        \end{itemize}
      \end{greencard}
    \end{column}
  \end{columns}
\end{frame}

\begin{frame}{Biblioteca e Ambientes de Aprendizagem}{\faBookReader\ Conhecimento Sem Limites: Acervo Físico, Digital e Espaços de Estudo}
  \begin{columns}[t]
    \begin{column}{0.48\textwidth}
      \begin{moderncard}{\faBook\ Acervo Físico e Sistema SophiA}
        \scriptsize
        \begin{itemize}\setlength{\itemsep}{0.25em}
          \item \textbf{Bibliografia Atualizada:} Títulos de referência em programação, arquitetura, dados, IA e redes.
          \item \textbf{Catálogo Informatizado (SophiA):} Consulta, reserva e renovação de empréstimos online.
          \item \textbf{Salas Climatizadas:} Espaços privativos para reuniões de equipe do PIM e estações individuais de estudo.
        \end{itemize}
      \end{moderncard}
    \end{column}
    
    \begin{column}{0.48\textwidth}
      \begin{purplecard}{\faGlobeAmericas\ Biblioteca Virtual \& Portal CAPES}
        \scriptsize
        \begin{itemize}\setlength{\itemsep}{0.25em}
          \item \textbf{Portal de Periódicos CAPES:} Acesso gratuito a publicações da IEEE, ACM e Springer.
          \item \textbf{Plataformas Virtuais:} Milhares de livros digitais disponíveis 24h por dia.
          \item \textbf{Suporte à Normalização:} Orientação técnica para trabalhos e artigos científicos na norma ABNT.
        \end{itemize}
      \end{purplecard}
    \end{column}
  \end{columns}
\end{frame}


% Secao 8: Como Entrar, Assistência Estudantil e Encerramento
% Secao 8: Como Entrar, Assistencia Estudantil e Encerramento

\begin{frame}{Como Ingressar em ADS no IFCE Tauá}{\faDoorOpen\ Seu Futuro Começa Aqui: 35 Vagas Anuais 100\% Gratuitas}
  \begin{columns}[t]
    \begin{column}{0.48\textwidth}
      \begin{greencard}{\faGraduationCap\ Formas de Ingresso}
        \begin{itemize}\setlength{\itemsep}{0.35em}\footnotesize
          \item \textbf{SISU / ENEM:} Principal porta de entrada utilizando a nota do Exame Nacional do Ensino Médio.
          \item \textbf{Vestibular Próprio IFCE:} Editais específicos divulgados periodicamente pelo instituto.
          \item \textbf{Vagas Remanescentes:} Editais para Transferência Externa e Portadores de Diploma de Curso Superior.
          \item \textbf{Periodicidade:} Entrada anual com turmas de 35 novos estudantes selecionados.
        \end{itemize}
      \end{greencard}
    \end{column}
    
    \begin{column}{0.48\textwidth}
      \begin{statcard}
        \Huge \textcolor{IFCEDarkGreen}{\textbf{35}} \\
        \small \textbf{Vagas Anuais Oferecidas} \\
        \scriptsize \textcolor{IFCERed}{\faCheckCircle\ Ensino Superior Federal 100\% Gratuito}
      \end{statcard}
      \vspace{0.2cm}
      \begin{tcolorbox}[enhanced, colback=IFCEGreen!8, colframe=IFCEGreen, arc=2mm, boxrule=0.8pt, left=2.5mm, right=2.5mm, top=2mm, bottom=2mm]
        \scriptsize \faStar\ \textbf{Sem mensalidades, sem taxas de matrícula:} Uma formação tecnológica de nível internacional mantida pelo Governo Federal.
      \end{tcolorbox}
    \end{column}
  \end{columns}
\end{frame}

\begin{frame}{Assistência Estudantil \& Permanência}{\faHands\ Suporte Completo para Você Ingressar, Permanecer e Concluir}
  \begin{columns}[t]
    \begin{column}{0.5\textwidth}
      \begin{bluecard}{\faHandHoldingUsd\ Auxílios Financeiros Diretos}
        \begin{itemize}\setlength{\itemsep}{0.25em}\scriptsize
          \item \textbf{Auxílio Moradia:} Para estudantes de outras cidades da região.
          \item \textbf{Auxílio Transporte \& Alimentação:} Apoio diário para o trajeto e refeições.
          \item \textbf{Auxílio Creche / Pais \& Mães Discentes:} Suporte aos estudantes com filhos.
          \item \textbf{Auxílio Didático-Pedagógico \& Óculos:} Recursos para aquisição de materiais e exames oftalmológicos.
          \item \textbf{Auxílio Pré-Embarque Internacional:} Custeio de documentos para o intercâmbio.
        \end{itemize}
      \end{bluecard}
    \end{column}
    
    \begin{column}{0.47\textwidth}
      \begin{purplecard}{\faHeartbeat\ Apoio Biopsicossocial \& NAPNE}
        \begin{itemize}\setlength{\itemsep}{0.25em}\scriptsize
          \item \textbf{Equipe Multidisciplinar:} Assistente Social, Psicólogo, Nutricionista, Enfermeira e Odontólogo no campus.
          \item \textbf{Inclusão com o NAPNE:} Núcleo especializado em acessibilidade, tecnologias assistivas e tutoria pedagógica adaptada.
          \item \textbf{Acolhimento Integral:} Cuidado permanente com a saúde mental e permanência acadêmica com êxito.
        \end{itemize}
      \end{purplecard}
    \end{column}
  \end{columns}
\end{frame}

\begin{frame}[plain]
  \begin{columns}[c]
    \begin{column}{0.38\textwidth}
      \centering
      \includegraphics[width=0.78\textwidth]{logo_ifce_taua.png}
      \vspace{0.3cm}
      
      {\tiny Rua Antônio Teixeira Benevides, 01 -- Colibris} \\
      \vspace{0.05cm}
      {\tiny CEP: 63660-000 \quad | \quad Tauá -- CE}
    \end{column}
    
    \begin{column}{0.60\textwidth}
      {\large \textbf{\textcolor{IFCEDarkGreen}{O Seu Futuro na Tecnologia Começa Agora!}}}
      
      \vspace{0.25cm}
      {\footnotesize \textbf{Venha transformar os Inhamuns e o mundo através do software.}}
      
      \vspace{0.35cm}
      \begin{moderncard}{\faInfoCircle\ Canais Oficiais de Contato}
        \begin{itemize}\setlength{\itemsep}{0.2em}\scriptsize
          \item \textbf{\faGlobe\ Site Oficial:} \textcolor{AccentBlue}{\textbf{www.ifce.edu.br/taua}}
          \item \textbf{\faEnvelope\ E-mail Institucional:} gabinete.taua@ifce.edu.br
          \item \textbf{\faPhone\ Telefone:} (88) 2134-1065
          \item \textbf{\faInstagram\ Redes Sociais:} @ifcetaua
        \end{itemize}
      \end{moderncard}
      
      \vspace{0.25cm}
      \centering
      \textbf{\textcolor{IFCERed}{\faRocket\ Esperamos por você no processo seletivo SISU / IFCE!}}
    \end{column}
  \end{columns}
\end{frame}


\end{document}
